\chapter{共享对象回收}
\label{ch:reclaim}
索引移除与读者退出可能发生在不同时间。$r_i$ 表示对象 $i$ 当前的读者引用计数。退役表示对象已移出索引，回收表示释放其存储空间。对象退役后仍可能被读者引用，因此仅凭退役状态无法判断对象是否可以回收。回收要求对象已经退役且 $r_i=0$，具体条件及检查与释放的互斥假设见算法说明。
\begin{quote}
\textbf{回收规则（教学伪代码）}\\
输入：对象退役状态 $retired_i$ 和读者引用计数 $r_i$。\\
若 $retired_i \land (r_i=0)$，释放对象存储；否则保留。\\
假设检查与释放由同一互斥区保护，移出索引后不能产生新引用。
\end{quote}

读者退出临界区后递减引用计数。
仅当对象已经退役且不再被读者引用时，回收器才释放对象空间。
\section{结果}
教学合成轨迹W中，延迟回收策略的峰值保留空间为72 KiB，基线为100 KiB，降低28\%。表~\ref{tab:space}给出两组轨迹的合成值。在轨迹V中，延迟回收策略的峰值保留空间为110 KiB，比基线100 KiB增加10\%。这两组教学合成轨迹表明，延迟回收策略的峰值保留空间随轨迹而异，轨迹W中的降低并未在轨迹V中出现。
\begin{table}[ht]
\centering
\caption{教学合成轨迹的峰值保留空间（KiB）；非系统实测。}\label{tab:space}
\begin{tabular}{lrr}
轨迹 & 基线 & 延迟回收策略 \\
W & 100 & 72 \\
V & 100 & 110
\end{tabular}
\end{table}

